\documentclass[11pt,a4paper]{article}

\usepackage[utf8]{inputenc}
\usepackage[T1]{fontenc}
\usepackage[english]{babel}
\usepackage{lmodern}
\usepackage[a4paper,margin=1in]{geometry}
\usepackage{setspace}
\usepackage{parskip}
\usepackage{ragged2e}
\usepackage{booktabs}
\usepackage{longtable}
\usepackage{tabularx}
\usepackage{array}
\usepackage{enumitem}
\usepackage{xcolor}
\usepackage{hyperref}
\usepackage{fancyhdr}
\usepackage{lastpage}

\hypersetup{
  colorlinks=true,
  linkcolor=blue!50!black,
  urlcolor=blue!50!black,
  pdftitle={Methodological Support Report v2 for the ANE Platform},
  pdfauthor={Codex},
  pdfsubject={Methodological analysis of the frontend and backend platform layers},
  pdfkeywords={software architecture, methodology, monitoring, frontend, backend, ANE}
}

\setstretch{1.12}
\setlist[itemize]{leftmargin=1.5em}
\setlist[enumerate]{leftmargin=1.5em}

\newcommand{\project}{ANE Spectral Sensing Platform}

\pagestyle{fancy}
\fancyhf{}
\lhead{\project}
\rhead{Methodological Support Report v2}
\cfoot{\thepage\ of \pageref{LastPage}}
\setlength{\headheight}{14pt}

\title{
  \vspace{-2em}
  \textbf{Methodological Support Report v2}\\
  \large Methods, Justifications, and Architectural Interpretation for \project
}
\author{Prepared from repository evidence and platform context}
\date{May 23, 2026}

\begin{document}
\pagenumbering{gobble}

\begin{titlepage}
  \centering
  \vspace*{3cm}
  {\LARGE \textbf{Methodological Support Report v2}\par}
  \vspace{0.6cm}
  {\Large Methods, Justifications, and Architectural Interpretation\par}
  \vspace{0.4cm}
  {\large \project\par}
  \vspace{2cm}
  \begin{tabular}{@{}p{0.30\textwidth}p{0.56\textwidth}@{}}
    Document type: & Technical-methodological support report \\
    Scope: & Interface and coordination layers of the platform \\
    Evidence date: & Repository snapshot analyzed on May 23, 2026 \\
    Analytical basis: & Context documents, diagrams, service descriptions, data evidence, static inspection \\
    Excluded from direct scope: & Planning artifacts, deep RF DSP internals, independent post-processing internals \\
  \end{tabular}
  \vfill
  {\large Prepared through architectural recovery, documentary triangulation, and methodological synthesis\par}
\end{titlepage}
\clearpage
\pagenumbering{arabic}

\begin{abstract}
This report provides a second, more methodological version of the architectural analysis for the ANE spectral sensing platform. Its purpose is to explain how the platform is organized, why its principal design choices are reasonable, and which operational rules support monitoring, campaigns, alerts, and compliance reporting, without centering the discussion on file structures or implementation paths. The analysis is based on contextual platform descriptions, frontend and backend flow diagrams, service and database documentation, and static inspection of the implemented solution. The document adopts an evidence-based architectural recovery approach: concerns are identified, system capabilities are grouped into viewpoints, and the resulting strategies are justified according to operational continuity, temporal coherence, maintainability, traceability, and security. The result is a report oriented toward methods, design rationale, and governance rather than code inventory.
\end{abstract}

\tableofcontents
\newpage

\section{Introduction}

\subsection{Problematic Introduction}
The platform studied in this repository addresses a complex operational problem: it must transform heterogeneous radio-frequency measurements produced by distributed field sensors into information that can be used for supervision, decision making, and compliance-oriented interpretation. This transformation is demanding because sensing activity is asynchronous, geographically distributed, and exposed to interruptions, stale data, environmental stress, and clock-related inconsistencies. At the same time, the human operators of the platform need both immediate situational awareness and delayed analytical workflows, including scheduled campaigns, alert review, and structured reporting.

Under these conditions, a merely descriptive inventory is insufficient. A methodological report must explain the operational logic of the system, the responsibilities assigned to the interface and coordination layers, the principles that justify their organization, and the procedural rules that give the platform predictable behavior over time. This second version is written with that objective in mind.

\subsection{Purpose}
The purpose of this document is to provide a methodological support report for the frontend and backend layers of \project. Specifically, it aims to:

\begin{itemize}
  \item define the operational problem that the platform solves;
  \item reconstruct the architectural logic of the current implementation from available evidence;
  \item explain the principal methods used to coordinate monitoring, campaigns, alerts, and reporting;
  \item justify the main design decisions according to quality attributes and domain constraints; and
  \item offer a reusable reference for technical reporting, governance, onboarding, and future redesign.
\end{itemize}

\subsection{Scope and Delimitation}
This report focuses on the platform layers that mediate between human operators, field sensors, stored information, and analytical services.

\begin{itemize}
  \item \textbf{Included}: access control, operational navigation, live monitoring, campaign management, alert handling, system configuration, service coordination, persistence strategy, temporal automation, and reporting orchestration.
  \item \textbf{Contextually considered}: the field sensing environment, the data-processing boundary, and the analytical post-processing stage. These are examined as dependencies and interfaces, not as the primary subject of the report.
  \item \textbf{Excluded}: planning-only materials, low-level signal-processing internals, and standalone analytical implementation details beyond what is necessary to explain platform behavior.
\end{itemize}

\section{Methodological Framework}

\subsection{Methodological Position}
This report adopts a descriptive and reconstructive software engineering methodology. It is descriptive because it documents the platform as it currently operates rather than as an idealized target design. It is reconstructive because the architecture is inferred from multiple forms of evidence, including domain descriptions, interaction diagrams, service documentation, data-model records, and static inspection of the implementation. The methodological stance is concern-oriented: the report is organized around the questions that matter for system understanding, namely access, monitoring, campaigns, alerts, reporting, persistence, temporal control, and maintainability.

\subsection{Evidence Base}
The analysis draws on five complementary evidence classes.

\begin{longtable}{@{}p{0.24\textwidth}p{0.26\textwidth}p{0.40\textwidth}@{}}
\toprule
\textbf{Evidence class} & \textbf{Observed content} & \textbf{Methodological contribution} \\
\midrule
\endhead
Platform context & General descriptions of mission, actors, workflows, and system responsibilities & Provides the domain frame that explains why the platform exists and what operational outcomes it must support. \\
Interaction and coordination diagrams & Narrative flows for the interface layer and the service layer & Clarify the sequence of decisions, the transfer of information, and the boundary between human actions and backend control. \\
Service and reporting descriptions & Descriptions of access, management, monitoring, campaigns, alerts, and report generation & Reveal the intended service taxonomy and the principal end-to-end use cases. \\
Database evidence & Tables, constraints, relationships, and time-series notes & Confirms the persistence model, identity anchors, temporal data strategy, and relational consistency assumptions. \\
Implementation inspection & Static review of representative frontend and backend behaviors & Verifies that the documented responsibilities correspond to actual operational mechanisms. \\
\bottomrule
\end{longtable}

\subsection{Analytical Procedure}
The analytical procedure followed four successive stages:

\begin{enumerate}
  \item \textbf{Documentary triangulation}: contextual descriptions, diagrams, and service summaries were compared to extract common vocabulary, system boundaries, and declared responsibilities.
  \item \textbf{Implementation verification}: the implemented frontend and backend behavior was inspected to confirm how those responsibilities are operationalized.
  \item \textbf{Architectural synthesis}: the recovered behaviors were grouped into coherent capabilities such as access control, device monitoring, campaign orchestration, alert fusion, reporting, and persistence governance.
  \item \textbf{Methodological justification}: the resulting design choices were interpreted through quality criteria including temporal coherence, maintainability, traceability, robustness, and security.
\end{enumerate}

\subsection{Evaluation Criteria}
Architectural decisions are interpreted through the following quality criteria:

\begin{itemize}
  \item \textbf{Functional adequacy}: whether the solution supports the platform's operational responsibilities.
  \item \textbf{Temporal coherence}: whether the design handles real-time change, scheduled processes, and stale-state risk in a controlled way.
  \item \textbf{Traceability}: whether it is possible to relate observable system behavior to identifiable platform capabilities.
  \item \textbf{Maintainability}: whether the decomposition supports controlled evolution and understandable reasoning.
  \item \textbf{Operational robustness}: whether the solution tolerates disconnections, partial data loss, and long-running workflows.
  \item \textbf{Security posture}: whether identity, authorization, and role boundaries are treated explicitly and consistently enough to support institutional use.
\end{itemize}

\section{System Foundation and Architectural Context}

\subsection{Domain Mission}
The domain mission of the platform is to convert raw spectrum observations into operationally meaningful knowledge. In practice, this means connecting people, remote sensors, stored histories, and analytical services so that the organization can observe the radio environment, respond to incidents, schedule measurement programs, and interpret findings under a compliance-oriented perspective.

\subsection{Operational Boundary}
The repository evidence reveals a four-part operational environment:

\begin{itemize}
  \item \textbf{Field sensing environment}: distributed devices that measure the radio spectrum, publish state and location, and receive measurement instructions.
  \item \textbf{Interface layer}: the user-facing environment through which operators authenticate, navigate, configure measurements, observe live behavior, manage campaigns, review alerts, and consult reports.
  \item \textbf{Coordination layer}: the service layer that stores information, enforces operational rules, synchronizes users and devices, and orchestrates reporting.
  \item \textbf{Analytical support environment}: external services or processes used for location enrichment and normative interpretation of collected measurements.
\end{itemize}

\subsection{Core Architectural Principle}
The dominant architectural principle is the separation between \textit{interaction} and \textit{authoritative control}. The interface layer organizes human workflows and presents data in an interpretable way, but the coordination layer owns the lasting operational state of sensors, campaigns, configurations, and reports. This principle is methodologically appropriate because monitoring systems must continue to behave coherently even when browsers disconnect, operators change screens, or multiple actors interact with the same sensing resources.

\subsection{Cross-Cutting Data Logic}
The data evidence points to three important structural facts:

\begin{itemize}
  \item device identity is anchored by a stable hardware identifier that acts as the practical reference across telemetry, configuration, and scheduling;
  \item campaign and reporting contexts combine normalized relational data with flexible configuration or result payloads, allowing both queryability and controlled variability; and
  \item time-dependent information is treated as a first-class concern, with explicit attention to historical telemetry, real-time status, and specialized storage patterns for time-series workloads.
\end{itemize}

\section{Frontend Methodological Analysis}

\subsection{Frontend Scope}
The interface layer is responsible for the human experience of the platform. Its mission is not limited to presentation; it structures the sequence through which a person enters the system, understands the current network situation, selects a sensor, initiates live observation, transitions from exploratory monitoring to scheduled campaigns, reviews incidents, and interprets results.

\subsection{Interaction Model}
The recovered interaction model follows a task-centered sequence:

\begin{enumerate}
  \item authenticate and validate session state before entering the operational workspace;
  \item obtain a general view of the network and device situation;
  \item navigate toward live monitoring, campaigns, alerts, or administration according to the task;
  \item configure and launch a measurement when direct observation is needed;
  \item observe live outputs such as spectrum, location, state, and audio when available;
  \item convert a useful measurement configuration into a reusable campaign when continuity is required; and
  \item interpret outcomes through historical review, alert examination, and report consultation.
\end{enumerate}

\subsection{Frontend Capability Map}
The interface layer can be understood through functional capabilities rather than implementation modules.

\begin{longtable}{@{}p{0.26\textwidth}p{0.22\textwidth}p{0.42\textwidth}@{}}
\toprule
\textbf{Capability} & \textbf{Primary concern} & \textbf{Methodological role} \\
\midrule
\endhead
Access and session control & Identity validation and protected entry & Ensures that operational screens are reached only after session state is coherent and permission checks are meaningful. \\
Operational overview & General situation awareness & Gives the user an initial mental model of network state, locations, and activity before more detailed actions are taken. \\
Device situational awareness & Sensor status, location, and association context & Supports informed sensor selection by exposing the condition and placement of field devices. \\
Live monitoring configuration & Immediate measurement preparation & Converts operational intent into a concrete acquisition setup that can be executed and later reused. \\
Spectral interpretation workspace & Real-time visualization and interaction & Makes dense technical measurements readable through charts, markers, units, and contextual controls. \\
Campaign management & Scheduled measurement planning and review & Bridges one-time observation and repeatable measurement programs governed by time. \\
Alert review & Incident awareness across current and past behavior & Prevents the operator from relying only on an instantaneous snapshot by preserving historical operational memory. \\
Administrative configuration & Resource and parameter governance & Keeps the platform adaptable by allowing managed change over devices, users, and general rules. \\
Audio-assisted observation & Supplemental live interpretation & Extends monitoring beyond visual traces when demodulated audio contributes to interpretation. \\
\bottomrule
\end{longtable}

\subsection{Frontend Design Choices and Their Justification}
\subsubsection{Authentication Before Operational Access}
The interface gives priority to session resolution before the user enters the main workspace. This choice is methodologically justified because it reduces ambiguity at the point where operational decisions begin. A monitoring interface should not oscillate between authenticated and unauthenticated states while the user is already evaluating sensors or measurements.

\subsubsection{Shared Operational Context Across Screens}
The platform maintains shared context around the currently selected device, measurement state, and active workflow. This is an appropriate design for a control-oriented environment because monitoring, configuration, alerts, and reporting are not isolated activities; they are different viewpoints over the same operational object.

\subsubsection{Hybrid Real-Time Interaction}
The monitoring experience combines event-style immediacy with periodic data refresh. This is methodologically coherent because not all payloads behave the same way. State changes and short notifications benefit from immediate delivery, whereas dense measurement arrays are better managed through controlled refresh cycles that protect interpretability and reduce synchronization complexity.

\subsubsection{Progressive Transition from Monitoring to Campaigns}
The interface allows a live configuration to become the seed of a scheduled activity. This is not a minor convenience; it is a method for preserving continuity between exploratory work and formal measurement plans. It reduces the cognitive and operational cost of translating a useful observation into a repeatable campaign.

\subsubsection{Alert Fusion Instead of Snapshot-Only Monitoring}
The platform does not rely only on the current state of a device. It combines present conditions with historical alert memory, which is methodologically stronger for operational supervision. A clean instantaneous reading does not erase earlier instability, and a historical alert without current context does not fully describe ongoing risk.

\subsection{Frontend Technical Families and Their Justification}
The interface layer relies on a family of architectural techniques whose value is methodological rather than merely technological.

\begin{longtable}{@{}p{0.26\textwidth}p{0.24\textwidth}p{0.40\textwidth}@{}}
\toprule
\textbf{Technical family} & \textbf{Observed role} & \textbf{Justification} \\
\midrule
\endhead
Component-based interface composition & Organization of screens and operational widgets & Supports complex workflows through reusable units without losing task-level clarity. \\
Typed interaction contracts & Explicit data structures between view logic and services & Reduces ambiguity in measurement, campaign, identity, and alert payloads. \\
Geospatial visualization & Device and campaign location awareness & Aligns naturally with field operations in which place is part of the interpretation. \\
High-density charting & Spectrum and historical trace interpretation & Makes measurement data readable at the scale required by technical supervision. \\
Stateful orchestration & Coordination of selected sensor, live acquisition, and workflow transitions & Preserves continuity across screens that share the same operational context. \\
Responsive operational layout & Support for dashboards, panels, and side-by-side inspection & Improves usability for long-lived sessions in which comparison and rapid navigation matter. \\
\bottomrule
\end{longtable}

\subsection{Frontend Procedural Rules and Their Justification}
Several frontend behaviors are better interpreted as explicit procedures rather than generic interface features.

\begin{longtable}{@{}p{0.28\textwidth}p{0.22\textwidth}p{0.40\textwidth}@{}}
\toprule
\textbf{Rule or procedure} & \textbf{Operational role} & \textbf{Justification} \\
\midrule
\endhead
Obsolescence control for live responses & Prevents delayed data from replacing the most recent user context & Necessary when operators switch devices or measurement ranges during active refresh cycles. \\
Trace reset and duplicate suppression & Keeps waterfall-style histories coherent after configuration changes & Prevents mixed histories and false visual continuity after the measured band changes. \\
Nearest-point interpretation for markers & Supports precise reading of discrete plotted samples & Appropriate when measurements are visualized as ordered bins rather than continuous analog curves. \\
Unit conversion for operator interpretation & Presents measurements in multiple meaningful forms & Helps technical users move between instrument-oriented and reporting-oriented views. \\
Local filtering and progressive review of campaign lists & Supports manageable inspection of moderate result sets & Improves usability without demanding unnecessary complexity in every query path. \\
Combined live-plus-historical issue evaluation & Merges present device condition with retained incident memory & Produces a more reliable operational picture than either mode on its own. \\
\bottomrule
\end{longtable}

\section{Backend Methodological Analysis}

\subsection{Backend Scope}
The coordination layer is the operational authority of the platform. Its responsibilities include validating access, receiving requests from the interface, accepting telemetry from sensors, storing state and history, coordinating live monitoring, managing campaigns, detecting alerts, communicating relevant updates in real time, and orchestrating report generation through external analytical support.

\subsection{Architectural Style}
The backend follows a concentrated coordination style rather than a highly fragmented service decomposition. This is methodologically appropriate for the domain because the hardest problem is not independent computation in many separate domains, but the consistent synchronization of users, sensors, measurements, schedules, and reports. A more distributed internal architecture would add communication overhead without necessarily improving the control of shared operational state.

\subsection{Backend Capability Map}
The coordination layer can be reconstructed through the following capabilities.

\begin{longtable}{@{}p{0.26\textwidth}p{0.22\textwidth}p{0.42\textwidth}@{}}
\toprule
\textbf{Capability} & \textbf{Primary concern} & \textbf{Methodological role} \\
\midrule
\endhead
Access and authorization governance & Identity, session, and role control & Separates institutional trust, platform session continuity, and privilege enforcement. \\
Resource administration & Management of sensors, antennas, users, and parameters & Maintains the controlled environment in which measurements can be executed reliably. \\
Monitoring orchestration & Live acquisition control and active configuration management & Ensures that a measurement is not just started, but also tracked, bounded, and recoverable. \\
Campaign orchestration & Scheduled measurement programs & Converts isolated observations into time-governed, repeatable operations. \\
Telemetry reception & Intake of device state, location, spectrum, and audio & Acts as the ingestion point through which the field environment becomes platform knowledge. \\
Persistence governance & Historical storage, current-state access, and data consistency & Preserves both transactional entities and temporal evidence required for reporting and audit. \\
Real-time dissemination & Immediate update propagation & Keeps operators informed without forcing every workflow into manual refresh. \\
Alert generation & Conversion of abnormal conditions into actionable records & Transforms raw telemetry into operationally meaningful incidents. \\
Reporting orchestration & Integration of measurements, location context, and analytical interpretation & Produces structured results that can support compliance-oriented reasoning. \\
Background automation & Periodic validation and lifecycle enforcement & Maintains temporal correctness even when no user is actively supervising a process. \\
\bottomrule
\end{longtable}

\subsection{Persistence Model and Data Governance}
The persistence strategy is hybrid. Stable domain entities such as people, devices, physical accessories, and campaigns are represented in relational form because they require integrity, identity, and controlled change. Variable configuration envelopes and analytical results use flexible structures because the domain includes evolving parameters and report payloads that do not fit a rigid schema equally well at all times. Historical telemetry is preserved as chronological evidence, which is consistent with a system where time, recurrence, and traceability are core concerns rather than secondary features.

This strategy is methodologically sound because the platform must serve two kinds of needs at once: durable transactional governance and high-volume temporal observation. The evidence also suggests a design under evolution, especially in the treatment of time-series information and timestamp representations. That is not a methodological flaw by itself, but it introduces an architectural obligation: all time-sensitive logic must explicitly account for representational differences so that schedules, freshness checks, and report windows remain consistent.

\subsection{Backend Design Choices and Their Justification}
\subsubsection{Backend as Source of Truth for Monitoring State}
The monitoring lifecycle is maintained as authoritative backend state rather than transient browser state. This is a strong architectural choice because sensing operations can outlive a single user session or screen. If the control layer did not own active monitoring state, the system would be vulnerable to accidental interruption whenever the interface disappears or connectivity degrades.

\subsubsection{Hybrid Identity Model}
The access design combines institutional identity assurance with a platform-specific authorization envelope. Methodologically, this balances two needs: trust in an external identity source and local control over roles, session lifetime, and application-specific permissions. It also reduces friction by aligning platform access with organizational identity practices.

\subsubsection{Capability-Oriented Service Separation}
The backend is organized around domain capabilities such as access, measurement control, campaigns, alerts, and reporting. This separation is justified because it preserves conceptual traceability. Even if some capabilities are operationally dense, the domain language remains visible, which is valuable for governance and future refactoring.

\subsubsection{Application-Level Temporal Automation}
The platform contains explicit time-driven rules for status decay, monitoring timeouts, and campaign transitions. This is methodologically appropriate because temporal behavior in this domain is not incidental; it is part of the business logic. Encoding those rules explicitly in the coordination layer makes them inspectable and auditable.

\subsubsection{Separated Audio Treatment}
Audio is treated as a distinct monitoring stream rather than as a trivial extension of generic event traffic. This is justified because audio has different framing, continuity, and listener expectations than scalar state updates. A specialized treatment reduces interference between ordinary control messages and higher-throughput interpretive media.

\subsubsection{Report Orchestration Through External Analytical Support}
The backend acts as the orchestrator of reporting rather than as the sole analytical engine. This is methodologically coherent because reporting in this domain depends on combining stored measurements with contextual enrichment and specialized analysis. A coordinating role allows the platform to integrate these steps while keeping the core platform centered on supervision and control.

\subsection{Backend Technical Families and Their Justification}
The backend employs several technical families that align with the domain problem.

\begin{longtable}{@{}p{0.26\textwidth}p{0.24\textwidth}p{0.40\textwidth}@{}}
\toprule
\textbf{Technical family} & \textbf{Observed role} & \textbf{Justification} \\
\midrule
\endhead
HTTP service coordination & Human-facing requests and configuration operations & Appropriate for administrative and consultative workflows that require explicit request-response semantics. \\
Event-based real-time messaging & Fast propagation of state changes and live updates & Reduces delay for operational awareness without overloading ordinary query channels. \\
Relational persistence with flexible payload support & Storage of governed entities plus variable configurations and reports & Balances integrity, queryability, and adaptability. \\
Time-series aware storage strategy & Historical telemetry capture at operational scale & Matches the need to preserve chronological evidence for monitoring and reporting. \\
Federated identity plus local session control & Access assurance and privilege enforcement & Supports institutional trust while retaining platform autonomy over authorization. \\
External analytical integration & Reporting enrichment and interpretation & Keeps specialized analysis available without dissolving the platform's coordinating role. \\
\bottomrule
\end{longtable}

\subsection{Backend Procedural Rules and Their Justification}
The backend contains deterministic operational rules that are central to platform reliability.

\begin{longtable}{@{}p{0.28\textwidth}p{0.22\textwidth}p{0.40\textwidth}@{}}
\toprule
\textbf{Rule or procedure} & \textbf{Operational role} & \textbf{Justification} \\
\midrule
\endhead
Sensor status decay state model & Classifies device availability according to recent activity & Necessary to distinguish active devices from delayed or disconnected ones in a meaningful way. \\
Threshold-based alerting with repetition control & Converts abnormal telemetry into actionable incidents without flooding operators & Supports vigilance while preserving the practical value of alert history. \\
Backend-enforced monitoring timeout & Stops stale or over-extended monitoring sessions & Protects resources and preserves operational discipline when interfaces disappear or sessions are abandoned. \\
Conflict detection for campaign assignment & Prevents overlapping scheduled use of the same sensing resource & Maintains exclusivity and avoids contradictory instructions to field devices. \\
Timezone-aware campaign transitions & Evaluates schedule changes in the intended operational timezone & Reduces ambiguity in date-time logic and aligns execution with the real operating calendar. \\
Configuration-to-device reconstruction rules & Converts stored measurement intent into device-usable parameters & Preserves compatibility when different configuration shapes must still lead to executable measurement behavior. \\
Threshold-sensitive cache reuse in reporting & Reuses prior analytical results only when the request is semantically compatible & Balances efficiency with correctness in report generation. \\
Priority-based location resolution & Chooses the best available location source for report context & Improves resilience when one source of geographic information is incomplete or missing. \\
Specialized audio frame handling & Preserves usable live audio for interpretation & Supports real-time auditory analysis without reducing the clarity of general platform messaging. \\
\bottomrule
\end{longtable}

\section{Cross-Cutting Architectural Interpretation}

\subsection{Observed Strengths}
From a methodological standpoint, the platform shows several strong characteristics.

\begin{itemize}
  \item \textbf{Clear operational decomposition}: the platform distinguishes access, monitoring, campaigns, alerts, administration, and reporting as recognizable concerns.
  \item \textbf{Strong control-plane orientation}: lasting operational state is governed centrally rather than being left to volatile user sessions.
  \item \textbf{Explicit temporal reasoning}: status windows, schedules, monitoring bounds, and historical traces are treated as design primitives.
  \item \textbf{Interpretability of operational rules}: many important behaviors are deterministic, auditable, and suitable for institutional supervision.
  \item \textbf{Continuity between exploration and formal measurement}: live observation can mature into scheduled campaigns and later into structured reporting.
\end{itemize}

\subsection{Observed Constraints and Design Tensions}
The evidence also reveals tensions that matter for future governance.

\begin{itemize}
  \item \textbf{Concentrated orchestration}: some responsibilities remain dense and centrally accumulated, which helps visibility but may increase long-term maintenance effort.
  \item \textbf{Mixed temporal representations}: differing timestamp conventions complicate reasoning across storage, scheduling, and reporting boundaries.
  \item \textbf{Evolutionary time-series strategy}: the persistence approach reflects ongoing refinement rather than a final stabilized model for all telemetry classes.
  \item \textbf{Security normalization pressure}: the access model is conceptually sound, but architectural maturity depends on making enforcement consistently uniform across all sensitive operations.
  \item \textbf{Dependence on external analytical availability}: reporting quality depends not only on stored data, but also on the availability of supporting enrichment and analysis services.
\end{itemize}

\section{Synthesis of Methods and Justifications}

\subsection{Why the Main Strategies Are Methodologically Coherent}
The core strategies of the platform align well with the problem it is solving. A backend-centered control plane is justified because devices, campaigns, and reports persist beyond any one user session. A hybrid real-time model is justified because not all monitoring payloads require the same transport semantics. Explicit temporal automation is justified because the domain depends on scheduled and freshness-sensitive behavior, not only on manual user action. Finally, report orchestration with contextual enrichment is justified because compliance-oriented interpretation demands more than raw measurement retrieval.

\subsection{Why the Technical Families Fit the Domain}
The selected technical families favor inspectability, controlled integration, and practical evolution. The interface layer uses architectural techniques suited to long-lived, stateful, information-dense workflows. The coordination layer uses patterns suited to governance, persistence, and orchestration. The data layer supports both structured administration and time-dependent evidence. Together, these choices form a pragmatic architecture for a supervision platform rather than an abstract demonstration system.

\subsection{Why the Operational Rules Matter}
The most important algorithms in the platform are not opaque prediction models; they are explicit operational procedures. Status decay, timeout enforcement, campaign conflict checks, alert deduplication, threshold-aware report reuse, and location prioritization are all understandable rules that can be inspected, audited, and revised. This is methodologically valuable in regulatory and supervisory contexts, where explainability is as important as raw automation.

\section{Conclusion}
The frontend and backend of \project form a coherent supervision platform whose principal architectural purpose is coordination. The interface layer structures operator activity, guiding users from access control to monitoring, campaigns, alerts, and reports. The backend consolidates authority over measurement state, persistence, timing, and analytical orchestration.

From a methodological perspective, the platform demonstrates a sound alignment between domain needs and architectural choices. Its main strengths are the clarity of its operational concerns, the central role of temporal control, and the use of auditable procedural rules. Its principal tensions involve concentrated orchestration, heterogeneous temporal representation, and the need for progressively stronger normalization of security and time-series strategy. Even so, the evidence supports the conclusion that the platform already constitutes a defensible technical foundation for spectral monitoring, campaign governance, alert supervision, and compliance-oriented interpretation.

This second version of the report therefore serves as a more abstract and method-focused basis for institutional reporting, architectural review, and future redesign because it explains not only what the platform contains, but why its current organization is methodologically reasonable.

\appendix

\section{Evidence Categories Considered}
\begin{itemize}
  \item General platform descriptions and mission statements.
  \item Frontend and backend interaction diagrams.
  \item Service and workflow documentation for access, monitoring, campaigns, alerts, and reports.
  \item Database inventory, relationships, constraints, and time-series notes.
  \item Static inspection of representative interface and coordination behaviors.
\end{itemize}

\section{Analytical Lenses Used}
\begin{itemize}
  \item Concern-oriented architectural recovery.
  \item Workflow reconstruction from user and service flows.
  \item Persistence and temporal-governance analysis.
  \item Design-rationale interpretation through quality attributes.
  \item Cross-cutting assessment of robustness, traceability, and maintainability.
\end{itemize}

\end{document}
